% Separate complete-legend versions; previous artwork is untouched.

\begin{figure}[!htbp]
    \centering
    \includegraphics[width=\textwidth]{SI_TiNi_2panel_side_colors_all_species.pdf}
    \caption{Pourbaix diagrams of TiNi. (a) \ce{NH3} and Gly; (b) \ce{NH3}, Gly, and \ce{CN-}. The metal-ion activity is $10^{-4}$. Ligand-containing panels use $[\ce{NH3}]_{\mathrm{tot}}=0.02\,\mathrm{M}$ and $[\mathrm{Gly}]_{\mathrm{tot}}=0.1\,\mathrm{M}$; cyanide-containing panels additionally use $[\ce{CN-}]_{\mathrm{tot}}=10^{-4}\,\mathrm{M}$. The black outline marks pH 10.5--13.5 and $-2.0$ to $2.0$ V vs RHE (converted to SHE and clipped to the plotted range).}
    \label{fig:Ti_Ni_Pourbaix_SI}
    \phantomsubcaption\label{fig:TiNi_Pourbaix_NH3_Gly_SI}
    \phantomsubcaption\label{fig:TiNi_Pourbaix_NH3_Gly_CN_SI}
\end{figure}

\begin{figure}[!htbp]
    \centering
    \includegraphics[width=\textwidth]{SI_TiCu_2panel_side_colors_all_species.pdf}
    \caption{Pourbaix diagrams of TiCu. (a) No added ligands; (b) \ce{NH3}, Gly, and \ce{CN-}. The metal-ion activity is $10^{-4}$. Ligand-containing panels use $[\ce{NH3}]_{\mathrm{tot}}=0.02\,\mathrm{M}$ and $[\mathrm{Gly}]_{\mathrm{tot}}=0.1\,\mathrm{M}$; cyanide-containing panels additionally use $[\ce{CN-}]_{\mathrm{tot}}=10^{-4}\,\mathrm{M}$. The black outline marks pH 10.5--13.5 and $-2.0$ to $2.0$ V vs RHE (converted to SHE and clipped to the plotted range).}
    \label{fig:Ti_Cu_Pourbaix_SI}
    \phantomsubcaption\label{fig:TiCu_Pourbaix_NH3_Gly_SI}
    \phantomsubcaption\label{fig:TiCu_Pourbaix_NH3_Gly_CN_SI}
\end{figure}

\begin{figure}[!htbp]
    \centering
    \includegraphics[width=\textwidth]{SI_TiFe_2panel_side_colors_all_species.pdf}
    \caption{Pourbaix diagrams of TiFe. (a) \ce{NH3} and Gly; (b) \ce{NH3}, Gly, and \ce{CN-}. The metal-ion activity is $10^{-4}$. Ligand-containing panels use $[\ce{NH3}]_{\mathrm{tot}}=0.02\,\mathrm{M}$ and $[\mathrm{Gly}]_{\mathrm{tot}}=0.1\,\mathrm{M}$; cyanide-containing panels additionally use $[\ce{CN-}]_{\mathrm{tot}}=10^{-4}\,\mathrm{M}$. The black outline marks pH 10.5--13.5 and $-2.0$ to $2.0$ V vs RHE (converted to SHE and clipped to the plotted range).}
    \label{fig:Ti_Fe_Pourbaix_SI}
    \phantomsubcaption\label{fig:TiFe_Pourbaix_NH3_Gly_SI}
    \phantomsubcaption\label{fig:TiFe_Pourbaix_NH3_Gly_CN_SI}
\end{figure}

\begin{figure}[!htbp]
    \centering
    \includegraphics[width=0.6\textwidth]{SI_TiMg_1panel_side_colors_all_species.pdf}
    \caption{Pourbaix diagrams of TiMg. (a) \ce{NH3} and Gly. The metal-ion activity is $10^{-4}$. Ligand-containing panels use $[\ce{NH3}]_{\mathrm{tot}}=0.02\,\mathrm{M}$ and $[\mathrm{Gly}]_{\mathrm{tot}}=0.1\,\mathrm{M}$; cyanide-containing panels additionally use $[\ce{CN-}]_{\mathrm{tot}}=10^{-4}\,\mathrm{M}$. The black outline marks pH 10.5--13.5 and $-2.0$ to $2.0$ V vs RHE (converted to SHE and clipped to the plotted range).}
    \label{fig:Ti_Mg_Pourbaix_SI}
    \phantomsubcaption\label{fig:TiMg_Pourbaix_NH3_Gly_SI}
\end{figure}

\begin{figure}[!htbp]
    \centering
    \includegraphics[width=0.6\textwidth]{SI_TiMn_1panel_side_colors_all_species.pdf}
    \caption{Pourbaix diagrams of TiMn. (a) \ce{NH3} and Gly. The metal-ion activity is $10^{-4}$. Ligand-containing panels use $[\ce{NH3}]_{\mathrm{tot}}=0.02\,\mathrm{M}$ and $[\mathrm{Gly}]_{\mathrm{tot}}=0.1\,\mathrm{M}$; cyanide-containing panels additionally use $[\ce{CN-}]_{\mathrm{tot}}=10^{-4}\,\mathrm{M}$. The black outline marks pH 10.5--13.5 and $-2.0$ to $2.0$ V vs RHE (converted to SHE and clipped to the plotted range).}
    \label{fig:Ti_Mn_Pourbaix_SI}
    \phantomsubcaption\label{fig:TiMn_Pourbaix_NH3_Gly_SI}
\end{figure}

\begin{figure}[!htbp]
    \centering
    \includegraphics[width=0.6\textwidth]{SI_TiZn_1panel_side_colors_all_species.pdf}
    \caption{Pourbaix diagrams of TiZn. (a) \ce{NH3}, Gly, and \ce{CN-}. The metal-ion activity is $10^{-4}$. Ligand-containing panels use $[\ce{NH3}]_{\mathrm{tot}}=0.02\,\mathrm{M}$ and $[\mathrm{Gly}]_{\mathrm{tot}}=0.1\,\mathrm{M}$; cyanide-containing panels additionally use $[\ce{CN-}]_{\mathrm{tot}}=10^{-4}\,\mathrm{M}$. The black outline marks pH 10.5--13.5 and $-2.0$ to $2.0$ V vs RHE (converted to SHE and clipped to the plotted range).}
    \label{fig:Ti_Zn_Pourbaix_SI}
    \phantomsubcaption\label{fig:TiZn_Pourbaix_NH3_Gly_CN_SI}
\end{figure}

\begin{figure}[!htbp]
    \centering
    \includegraphics[width=\textwidth]{SI_PdAg_2panel_side_colors_all_species.pdf}
    \caption{Pourbaix diagrams of PdAg. (a) \ce{NH3} and Gly; (b) \ce{NH3}, Gly, and \ce{CN-}. The metal-ion activity is $10^{-4}$. Ligand-containing panels use $[\ce{NH3}]_{\mathrm{tot}}=0.02\,\mathrm{M}$ and $[\mathrm{Gly}]_{\mathrm{tot}}=0.1\,\mathrm{M}$; cyanide-containing panels additionally use $[\ce{CN-}]_{\mathrm{tot}}=10^{-4}\,\mathrm{M}$. The black outline marks pH 10.5--13.5 and $-2.0$ to $2.0$ V vs RHE (converted to SHE and clipped to the plotted range).}
    \label{fig:Pd_Ag_Pourbaix_SI}
    \phantomsubcaption\label{fig:PdAg_Pourbaix_NH3_Gly_SI}
    \phantomsubcaption\label{fig:PdAg_Pourbaix_NH3_Gly_CN_SI}
\end{figure}

\begin{figure}[!htbp]
    \centering
    \includegraphics[width=\textwidth]{SI_PdCu_2panel_side_colors_all_species.pdf}
    \caption{Pourbaix diagrams of PdCu. (a) \ce{NH3} and Gly; (b) \ce{NH3}, Gly, and \ce{CN-}. The metal-ion activity is $10^{-4}$. Ligand-containing panels use $[\ce{NH3}]_{\mathrm{tot}}=0.02\,\mathrm{M}$ and $[\mathrm{Gly}]_{\mathrm{tot}}=0.1\,\mathrm{M}$; cyanide-containing panels additionally use $[\ce{CN-}]_{\mathrm{tot}}=10^{-4}\,\mathrm{M}$. The black outline marks pH 10.5--13.5 and $-2.0$ to $2.0$ V vs RHE (converted to SHE and clipped to the plotted range).}
    \label{fig:Pd_Cu_Pourbaix_SI}
    \phantomsubcaption\label{fig:PdCu_Pourbaix_NH3_Gly_SI}
    \phantomsubcaption\label{fig:PdCu_Pourbaix_NH3_Gly_CN_SI}
\end{figure}

\begin{figure}[!htbp]
    \centering
    \includegraphics[width=\textwidth]{SI_PdNi_2panel_side_colors_all_species.pdf}
    \caption{Pourbaix diagrams of PdNi. (a) \ce{NH3} and Gly; (b) \ce{NH3}, Gly, and \ce{CN-}. The metal-ion activity is $10^{-4}$. Ligand-containing panels use $[\ce{NH3}]_{\mathrm{tot}}=0.02\,\mathrm{M}$ and $[\mathrm{Gly}]_{\mathrm{tot}}=0.1\,\mathrm{M}$; cyanide-containing panels additionally use $[\ce{CN-}]_{\mathrm{tot}}=10^{-4}\,\mathrm{M}$. The black outline marks pH 10.5--13.5 and $-2.0$ to $2.0$ V vs RHE (converted to SHE and clipped to the plotted range).}
    \label{fig:Pd_Ni_Pourbaix_SI}
    \phantomsubcaption\label{fig:PdNi_Pourbaix_NH3_Gly_SI}
    \phantomsubcaption\label{fig:PdNi_Pourbaix_NH3_Gly_CN_SI}
\end{figure}

\begin{figure}[!htbp]
    \centering
    \includegraphics[width=\textwidth]{SI_PdPt_2panel_side_colors_all_species.pdf}
    \caption{Pourbaix diagrams of PdPt. (a) \ce{NH3} and Gly; (b) \ce{NH3}, Gly, and \ce{CN-}. The metal-ion activity is $10^{-4}$. Ligand-containing panels use $[\ce{NH3}]_{\mathrm{tot}}=0.02\,\mathrm{M}$ and $[\mathrm{Gly}]_{\mathrm{tot}}=0.1\,\mathrm{M}$; cyanide-containing panels additionally use $[\ce{CN-}]_{\mathrm{tot}}=10^{-4}\,\mathrm{M}$. The black outline marks pH 10.5--13.5 and $-2.0$ to $2.0$ V vs RHE (converted to SHE and clipped to the plotted range).}
    \label{fig:Pd_Pt_Pourbaix_SI}
    \phantomsubcaption\label{fig:PdPt_Pourbaix_NH3_Gly_SI}
    \phantomsubcaption\label{fig:PdPt_Pourbaix_NH3_Gly_CN_SI}
\end{figure}

\begin{figure}[!htbp]
    \centering
    \includegraphics[width=\textwidth]{SI_AuPd_2panel_side_colors_all_species.pdf}
    \caption{Pourbaix diagrams of AuPd. (a) No added ligands; (b) \ce{NH3}, Gly, and \ce{CN-}. The metal-ion activity is $10^{-4}$. Ligand-containing panels use $[\ce{NH3}]_{\mathrm{tot}}=0.02\,\mathrm{M}$ and $[\mathrm{Gly}]_{\mathrm{tot}}=0.1\,\mathrm{M}$; cyanide-containing panels additionally use $[\ce{CN-}]_{\mathrm{tot}}=10^{-4}\,\mathrm{M}$. The black outline marks pH 10.5--13.5 and $-2.0$ to $2.0$ V vs RHE (converted to SHE and clipped to the plotted range).}
    \label{fig:Au_Pd_Pourbaix_SI}
    \phantomsubcaption\label{fig:AuPd_Pourbaix_NH3_Gly_SI}
    \phantomsubcaption\label{fig:AuPd_Pourbaix_NH3_Gly_CN_SI}
\end{figure}
